\section{Axiom Inheritance and the Non-Self-Balancing Finding}
\label{sec:axiom_inheritance}

Since $\volRlower$ is a lower bound on $\volR$, and $\volR =
\volL$ restricted to the exercised sub-poset, the axiom
inheritance analysis applies directly, with an additional layer:
the lower bound itself is not a measure (it is a bound on a
measure).

\begin{theorem}[Axiom-Analogue Analysis for $\volRlower$]
\label{thm:axiom_inheritance}
$\volRlower$ is a lower-bound statistic, not a poset measure in
the sense of~\citep[Proposition~1]{lasser2026poset}; the clauses
below describe \emph{operational analogues} of M1--M6 for
$\volRlower$, not literal axiom inheritance.  M4 and the
structural M6-failure result are substantive; M1, M2, M3, and
M5 are weaker conditional statements noted for completeness.

\begin{description}
\item[(M1) Non-negativity (analogue).] Satisfied.  $\volRlower$ is
  a sum of non-negative terms ($\Delta\volL(X_n) \geq 0$ per event;
  max-attribution preserves non-negativity).

\item[(M2) Null empty set (analogue).] Satisfied.  With zero
  sacrifice events, $\volRlower = 0$.

\item[(M3) Monotonicity (conditional analogue).] Satisfied in the
  observational direction: adding sacrifice events covering a
  larger subset of $U$ weakly increases $\volRlower$
  (Proposition~\ref{prop:monotone}).  M3 does \emph{not} hold
  for the relationship between $\volL$ and $\volRlower$: adding
  a benchmarked capability to the poset increases $\volL$ but
  does not affect $\volRlower$ unless that capability is
  subsequently sacrificed.  This is a lower-bound statistic
  property, not a poset-measure axiom.

\item[(M4) Additivity under disjointness (substantive).]
  Satisfied.  By construction, $\volR$ is $\volL$ restricted
  to the exercised sub-poset $P^{\mathrm{ex}} \subseteq P$.  M4
  for $\volL$ on $P$ states that $\volL$ is additive over
  components of $P$ that are pairwise poset-disjoint (no
  subsumption or cooperative-composition links).  Let $A, B
  \subseteq P^{\mathrm{ex}}$ be poset-disjoint in $P^{\mathrm{ex}}$.
  Since $P^{\mathrm{ex}}$ inherits the relation structure of
  $P$, $A$ and $B$ are also poset-disjoint in $P$.  M4 on $P$
  gives $\volL(A \cup B) = \volL(A) + \volL(B)$; restricting
  both sides to $P^{\mathrm{ex}}$ gives $\volR(A \cup B) =
  \volR(A) + \volR(B)$, establishing M4 for $\volR$.
  This inheritance is load-bearing for the aggregate bound
  (Theorem~\ref{thm:aggregate_bound}, Part~A).

\item[(M5) Non-triviality (conditional analogue).] A single
  sacrifice event with $\Delta\volL(X) > 0$ produces
  $\volRlower > 0$ on the capabilities the event covers.  This
  is an observation-dependent analogue, not the structural
  non-triviality of every valid capability.

\item[(M6) Superadditivity under independence (substantive).]
  Does \emph{not} unconditionally hold.  Two distinct failures,
  each witnessed by a constructive counterexample
  (Examples~\ref{ex:m6_fail_volR} and~\ref{ex:m6_fail_volRlower}
  below):
  \begin{enumerate}
  \item[(a)] \textbf{$\volR$ itself fails M6 (structural).}
    $\volR = \volL$ restricted to the exercised sub-poset, and
    merging two groups can change exercise indicators (one group
    provides alternative pathways for the other's outputs, making
    previously load-bearing capabilities dormant).  This failure
    is about the \emph{measure} $\volR$, independent of the
    observation mechanism.
  \item[(b)] \textbf{$\volRlower$ fails M6 (observational).}
    Merging two groups of sacrifice events changes per-capability
    attributions under max-attribution: the joint maximum over
    all events is less than the sum of within-group maxima.  This
    failure is about the \emph{lower bound}, a weaker claim---a
    lower bound failing superadditivity does not imply the
    underlying quantity fails superadditivity.
  \end{enumerate}
\end{description}
\end{theorem}

\begin{example}[$\volR$ M6-failure via alternative pathways]
\label{ex:m6_fail_volR}
Consider a single-agent setting with two poset-independent
capability groups:
\begin{itemize}
\item $A = \{c_A\}$ where $c_A = $ ``drive to gas station
  $\alpha$'', with weight $w(c_A) = 1$.
\item $B = \{c_B\}$ where $c_B = $ ``drive to gas station
  $\beta$'' (different location), with weight $w(c_B) = 1$.
\end{itemize}
$A$ and $B$ share no subsumption or cooperative-composition
relation: they target different fuel-access capabilities with no
shared participants.  Apply the exercise indicator
$e_t(d)$ of Appendix~\ref{app:exercise_indicator}: $d$ is
load-bearing at time $t$ iff removing $d$ would make some
realized cooperative output unrealizable.

In group $A$ alone, the agent uses $c_A$ for fuel access;
removing $c_A$ blocks fuel access, so $c_A$ is load-bearing.
$\volR(A) = w(c_A) = 1$.

In group $B$ alone, $c_B$ is similarly load-bearing.
$\volR(B) = 1$.

In the merged group $A \cup B$, the agent uses the closer of the
two (say $c_A$).  Removing $c_B$ does not block fuel access:
$c_A$ remains available.  Therefore $c_B$ is no longer
load-bearing.  $\volR(A \cup B) = w(c_A) = 1$.

M6 (superadditivity under poset-independence) predicts
$\volR(A \cup B) \geq \volR(A) + \volR(B) = 2$, but
$\volR(A \cup B) = 1 < 2$.  M6 fails for $\volR$ in the
subadditive direction---merging poset-independent groups can
reduce the count of load-bearing capabilities because alternative
pathways become available.
\end{example}

\begin{example}[$\volRlower$ M6-failure via max-attribution]
\label{ex:m6_fail_volRlower}
Consider two sacrifice events with the same bundle
$Y = \{a, b\}$ but different decomposition coefficients
(permitted under S5, since event-local decomposition is fixed at
observation time):
\begin{itemize}
\item Event $E_1$: $\Delta\volL(X_1) = 1$, $\alpha_{1,a} = 0.5$,
  $\alpha_{1,b} = 0.5$.
\item Event $E_2$: $\Delta\volL(X_2) = 1$, $\alpha_{2,a} = 0.3$,
  $\alpha_{2,b} = 0.7$.
\end{itemize}
Let group $A = \{E_1\}$ and group $B = \{E_2\}$.  Under
Theorem~\ref{thm:aggregate_bound} Part~B:
\begin{align*}
  \volRlower(a \mid A) &= 0.5, &
  \volRlower(b \mid A) &= 0.5, \\
  \volRlower(a \mid B) &= 0.3, &
  \volRlower(b \mid B) &= 0.7.
\end{align*}
Summing per-capability:
$\volRlower(A) = 1.0$, $\volRlower(B) = 1.0$.

In the merged group $A \cup B$, max-attribution over both events:
$\volRlower(a \mid A \cup B) = \max(0.5, 0.3) = 0.5$ and
$\volRlower(b \mid A \cup B) = \max(0.5, 0.7) = 0.7$.
Total: $\volRlower(A \cup B) = 1.2$.

M6 (superadditivity) predicts $\volRlower(A \cup B) \geq
\volRlower(A) + \volRlower(B) = 2.0$, but $\volRlower(A \cup B)
= 1.2 < 2.0$.  M6 fails for $\volRlower$ in the subadditive
direction---max-attribution over merged events takes the larger
single contribution per capability rather than summing
across groups, so the joint aggregate is bounded above by the
sum of within-group aggregates.

This is a property of the \emph{lower-bound construction}, not of
the underlying $\volR$: the true $\volR$ aggregate of the merged
events need not be subadditive, but our observational lower
bound is.  The non-self-balancing corollary below therefore
relies on failure~(a), not failure~(b).
\end{example}

\begin{corollary}[$\volR$ Is Not Self-Balancing]
\label{cor:not_self_balancing}
The self-balancing property does \emph{not} transfer from $\volL$
to $\volR$.  This follows from failure~(a) above: $\volR$ as a
measure does not satisfy M6 unconditionally, and M6 is
load-bearing for the self-balancing theorem
of~\citep[Theorem~1]{lasser2026poset}.

The corollary does \emph{not} follow from $\volRlower$ failing M6
alone---a lower bound failing measure axioms is unremarkable.
The corollary requires the structural argument from~(a), which
holds because merging groups can change exercise status.
\end{corollary}

This is the structural reason to use $\volR$ as a diagnostic
rather than an objective.  An actor maximizing $\volR$ would lack
the automatic diversity-preservation that makes $\volL$ safe.

\begin{remark}[When M6 does hold for $\volRlower$]
M6 holds for $\volRlower$ when the merged groups have sacrifice
events in poset-independent bundle categories
(Definition~\ref{def:category_partition})---each unbenchmarked
capability appears in sacrifice events from only one group, and
cross-group capabilities are poset-independent.  Under this
non-redundancy condition, $\volRlower$ inherits M1--M6.  Even
so, the self-balancing property transfers to $\volRlower$ only as
a \emph{lower bound} on self-balancing---$\volRlower$ being
self-balancing does not imply $\volR$ itself is self-balancing.
\end{remark}
